\section*{Part III: Heterogeneity and the Shape of Demand}

\subsection*{17. There Is No Average Customer Who Actually Buys}

A founder asks 40 prospects how much they would pay for the same workshop and averages the replies. The number has mathematical meaning; it may have almost no commercial meaning. Some imagine a recorded lecture, some a private critique, some a credential, some a deadline-driven engineering intervention. Their jobs and expectations differ. Averaging those people produces a fictional customer whose preferences resemble no 1 in the sample.

Customer heterogeneity has at least 6 sources: different desired outcomes; different severity and urgency; different alternatives; different beliefs about delivery; different financial and time constraints; and different decision rights. Each produces not merely a different number but sometimes a different product. A serious firm therefore estimates demand conditional on a precise segment and a precise offer.

For Bayes, ``people aged 21 to 30 interested in AI'' is unlikely to explain price acceptance well. ``Working backend engineers already assigned a consequential AI implementation, without experienced internal reviewers, able to allocate a specified number of hours'' may be more useful. The 2nd group could still reject the offer, but at least the hypothesis describes a coherent purchasing situation.

\subsection*{18. Segmentation by Circumstance Beats Segmentation by Convenient Labels}

Age, city, salary, occupation, and company size can help find buyers. They are often weak proxies for the actual reason the buyer purchases. A better segmentation logic begins with the \emph{job}, \emph{trigger}, \emph{current alternative}, \emph{consequences of failure}, and \emph{ability to fund} (Christensen et al. 2016). These conditions predict what a prospect considers valuable.

Consider 3 customers in the same engineering job title. 1 needs foundational concepts and may buy a low-cost self-paced course. Another has an upcoming project review and wants serious architecture feedback. A 3rd is already strong but wants access to research collaboration. A single 6-month curriculum with 1 price is unlikely to be equally compelling to all 3.

Build a segment dossier that includes observed past purchases, buyer decision context, desired transformation, current workaround, real spending limits, time capacity, credibility requirements, and the strongest reason to say no. Segment attractiveness belongs in Chapters 5 to 7; here the narrower task is to model the \emph{heterogeneity of reservation values and purchase feasibility} within that segment.

\subsection*{19. A Demand Curve Is a Distribution of Choices Under Assumptions}

Suppose a hypothetical population contains 100 otherwise comparable qualified buyers. At ₹10,000, 60 might buy; at ₹20,000, 35; at ₹30,000, 18. This gives a downward-sloping \emph{illustrative} relation between price and quantity. It is not evidence that a real group will behave so cleanly. Actual demand can have kinks, price-as-quality effects over limited ranges, heterogeneous buyer mixes, seasonality, and word-of-mouth dynamics.

\begin{equation}
Q(P) = N \cdot \Pr(\text{WTP}_i \ge P \mid \text{offer, context, awareness}).
\end{equation}

The expression assumes buyers know and can access the offer, and compresses budgeting and decision frictions into the reservation construct. More useful in practice is the full funnel: eligible audience $\rightarrow$ reached $\rightarrow$ understands $\rightarrow$ qualified $\rightarrow$ believes $\rightarrow$ attempts purchase $\rightarrow$ pays $\rightarrow$ begins $\rightarrow$ obtains value. A price change can affect multiple transitions. The estimated willingness-to-pay distribution of warm alumni referrals cannot be applied mechanically to cold prospects acquired through advertising.

A founder should resist treating 1 visible conversion percentage as a stable demand function. The question is not just ``What fraction bought?'' but ``Which people, through which route, with which alternatives, and at what stage of product maturity?''

\subsection*{20. Elasticity Measures Response; It Does Not Reveal Motive}

Price elasticity describes the responsiveness of demand to price. In differential form:

\begin{equation}
\epsilon = \left( \frac{dQ}{dP} \right) \left( \frac{P}{Q} \right).
\end{equation}
In ordinary downward-sloping settings, it is negative. A segment with many close substitutes often exhibits greater sensitivity than a segment facing a costly, urgent problem. But low measured sensitivity can also reflect switching friction, inertia, procurement lock-in, or a small experimental price range.

A price test may show weak response between ₹20,000 and ₹22,000 and a dramatic collapse at ₹25,000, where the buyer crosses a budgeting or approval threshold. 1 slope will not describe both regions. Response also changes with reputation, available alternatives, the stage of the purchasing journey, and customer experience.

Do not infer that a buyer who is ``not price sensitive'' does not care about money. They may care profoundly about avoiding failure, and the purchase may be modest relative to the potential downside. Elasticity is an observed property of choices, not a psychological diagnosis.
